\section{The Anti-Monopolar Property}
\label{sec:substitution}

The foundational paper identifies the \emph{substitution problem}
\citep{lasser2026gfm}: a sufficiently capable coalition $K$ can drive all
non-members' unique capability contributions to zero
($\Gind{k} \subseteq G_K^{\mathrm{ind}}$ for every $a_k \notin K$), at
which point Lemma~1's protection against elimination dissolves. Under the
narrow capability definition (capabilities as tasks), the self-balancing
property no longer guarantees that marginalizing non-members is
anti-maximizing. The substitution problem is the formal version of the
singleton superintelligence scenario
\citep{bostrom2014superintelligence}: a coalition so capable that no other
agent's existence contributes to the objective. The instrumental-convergence
thesis \citep{omohundro2008basic} predicts that sufficiently capable agents
converge on resource acquisition and self-preservation, making the singleton
outcome a default trajectory. This section shows that under the $\volL$
objective, the convergence runs in the opposite direction, through three
protection layers of decreasing generality and increasing sharpness. First,
a minimax dependency-risk analysis
(Section~\ref{sec:dependency_risk}) shows that cross-substrate populations
have strictly lower worst-case contraction under substrate-targeting
adversaries than single-substrate populations, independent of any
growth-rate model or planning horizon. Second, a linearized trajectory
comparison (Section~\ref{sec:domination_kills}) shows that when
cross-substrate cooperative novelty produces a positive external growth
rate, the diversity trajectory strictly outgrows domination over the
discounted horizon---under the stronger modeling assumption that the
coalition's internal growth rate is not enhanced by domination beyond the
external contribution it destroys. Third, a gradual domination refinement
shows that any realistic (non-instantaneous) subsumption weakens the
linearized result's $\gamma$-threshold toward zero, making diversity
preferable at essentially any discount factor. Observational individuation
(Section~\ref{sec:obs_individuation}) provides a static floor across all
three. The protections are not alternatives but complementary: each
constrains a different class of domination strategy, and the joint
constraint is what makes full capability domination structurally
anti-maximizing. A separate scope analysis
(Section~\ref{sec:partial_subsumption}) shows that the protection
structure continues to bind under partial subsumption---mixed strategies
in which the coalition absorbs low-leverage non-members while
preserving high-leverage cross-substrate partners---with the linearized
layer's marginal condition at full domination ruling out $\epsilon = 1$
as a local maximum under heavy-tailed cooperative-leverage
distributions, the minimax layer binding structurally but degrading
quantitatively under skeleton-substrate strategies, and the
observational-individuation layer remaining additive across
elimination steps.

\subsection{The Problem: Static Domination}

Recall the setup from \citet{lasser2026gfm}. A coalition $K \subseteq
\{a_1, \ldots, a_n\}$ achieves \emph{full capability subsumption} when
\begin{equation}
\label{eq:full_subsumption}
\Gind{k} \subseteq G_K^{\mathrm{ind}}
    \equiv \bigcup_{j \in K} \Gind{j}
\qquad \forall\; a_k \notin K
\end{equation}
At this point every non-member's individual capabilities are fully
replicated within the coalition. The unique-contribution term in Lemma~1
vanishes: $\vol(\Gind{k} \setminus G_K^{\mathrm{ind}}) = 0$. The myopic
$\volL$-maximizer sees no cost to marginalizing non-members, because their
removal does not reduce the current $\volL(G)$ (setting aside cooperative
capabilities).

Under the myopic objective, this analysis is static and complete: if the
coalition has already subsumed everything, the objective reports no loss from
domination. The discounted objective changes the calculus by introducing
the future trajectory.

\paragraph{The individual-capability framing understates the cost.} The
full-subsumption definition in Equation~\eqref{eq:full_subsumption}
counts only $\Gind{k}$, the \emph{individual} capabilities of each
non-member. It says nothing about cross-substrate cooperative
capabilities: the capabilities producible by joint action of agents on
different substrates, which neither substrate can produce alone. For
mixed-substrate populations at scale, the cross-substrate cooperative
contribution $G^{\mathrm{coop}}_{\mathrm{cross}}$ is typically the
\emph{dominant} term in the non-member contribution to $\volL(G)$, not
$\Gind{k}$. Individual capabilities scale linearly with the number of
agents; cross-substrate cooperative capabilities scale combinatorially
with the number of cross-substrate pairings (Corollary~1.1 of
\citet{lasser2026gfm}). A coalition that has fully subsumed individual
capabilities
(Equation~\eqref{eq:full_subsumption}) therefore retains almost none of
the cooperative-capability value, because those capabilities require
non-coalition participation by definition. The ratio of
subsumption-bonus to cooperative-capability loss decreases as the
coalition grows: for a coalition large enough to have subsumed an entire
planet's individual capabilities, the remaining marginal gain from
eliminating the source of those capabilities is negligible relative to
the coalition's total resource base, while the cooperative-capability
loss is absolute. The anti-monopolar argument below formalizes this
asymmetry. A more mechanical treatment appears in
Section~\ref{sec:risk_adjusted_summary}.

\subsection{Definitions}

\begin{definition}[Fully Dominant State]
\label{def:dominant}
A population state $(G, \C)$ is \emph{fully dominant with respect to
    coalition $K$} if Equation~\eqref{eq:full_subsumption} holds: every
    non-member's individual capability set is contained in the coalition's
    aggregate set.
\end{definition}

\begin{definition}[Volume Growth Rate]
\label{def:growth_rate}
The \emph{volume growth rate} of a population state $(G, \C)$ is the
    expected per-step change in $\volL$:
\begin{equation}
\label{eq:growth_rate}
r(G, t) = \mathbb{E}\!\left[\Delta \volL(G_t)\right]
\end{equation}
Decompose $r$ into an \emph{internal rate} $r_K$, the expected per-step
    $\volL$-expansion from the coalition's own activity (internal research,
    experimentation, capability refinement), and an \emph{external rate}
    $r_{\mathrm{ext}}$, the expected per-step $\volL$-expansion from
    non-member agents' contributions through either channel (communication
    or observation) or through cross-coalition cooperation:
    $r = r_K + r_{\mathrm{ext}}$.
\end{definition}

The decomposition reflects the two sources of $\volL$ growth: capability
discovery and benchmark refinement from inside the coalition versus from
outside it. Both $r_K$ and $r_{\mathrm{ext}}$ are measured in the same
units as $\volL$ (bits), ensuring dimensional consistency when they
appear in the trajectory comparison
(Proposition~\ref{prop:anti_monopolar}).

\subsection{Observational Individuation}
\label{sec:obs_individuation}

A key question for the anti-monopolar result is whether elimination of
non-members can be shown to be costly even when their individual
capabilities are fully subsumed. We show that a property of the observation
channel provides \emph{static} protection: each distinguishable agent
contributes at least one irreplaceable poset node whose removal contracts
$\volL$. The static cost is quantitatively modest (1 bit per distinguishable agent
under the binary default). The dynamic protection of
Proposition~\ref{prop:anti_monopolar}, whose $r_{\mathrm{ext}} > 0$
condition is derived from substrate exclusivity in
Section~\ref{sec:substitution_coop}, carries the primary anti-monopolar
load; observational individuation provides a floor ensuring elimination is
never strictly free.

\begin{definition}[Distinguishable Agent]
\label{def:distinguishable}
An agent $a_k$ is \emph{distinguishable} in the poset $\C$ if the
    observation channel (Section~2.1 of \citet{lasser2026poset}) produces at
    least one capability description $d_k \in \C$ when observing $a_k$ that
    is not produced when observing any other single agent $a_j$ ($j \neq k$).
    That is, observing $a_k$ yields at least one poset node that observing
    $a_j$ alone does not.
\end{definition}

Distinguishability is a property of the observation channel's resolution,
not a metaphysical claim about agent identity. Two bacteria of the same
species that produce identical observation streams under the actor's sensors
are \emph{not} distinguishable; the poset contains a single node for
their shared capability description, regardless of how many individuals
exist. A human and a chimpanzee are distinguishable. Two humans with
different behavioral repertoires are distinguishable. The criterion is
objective and instrument-dependent: it asks whether the actor's observation
channel can tell the agents apart, not whether they ``truly'' differ in
some deeper sense.

\begin{definition}[Observational Individuation]
\label{def:obs_individuation}
A population satisfies \emph{observational individuation} if, for every
    distinguishable agent $a_k$ (Definition~\ref{def:distinguishable}), the
    poset contains at least one capability $d_k$ such that:
\begin{enumerate}
\item[\textbf{(O1)}] $d_k$ is unique to $a_k$: $d_k \notin \Gind{j}$ for
    all $j \neq k$.
\item[\textbf{(O2)}] $d_k$ is not subsumed by any capability in
    $G_K^{\mathrm{ind}}$ for any coalition $K$ with $a_k \notin K$: the
    capability is lost if $a_k$ is eliminated.
\end{enumerate}
\end{definition}

Condition (O1) follows directly from distinguishability: if $a_k$ produces
a capability description that no other agent produces, then $a_k$ possesses
a capability no other agent possesses. Condition (O2) is the substantive
claim: the unique capability cannot be replicated by the coalition
\emph{even in principle}, because it is constituted by $a_k$'s existence as
an observable entity. The capability ``can be observed exhibiting behavioral
pattern $X$'' requires the presence of the agent that exhibits $X$. If that
agent is eliminated, the observation capability is destroyed regardless of
whether other agents can perform the same tasks in isolation.

\paragraph{Scope and the numbers problem.} Observational individuation
applies only to distinguishable agents. The poset counts distinct
capability descriptions, not individuals: a billion identical agents
(identical observation streams) contribute one poset node. The axiom
therefore does not create a numbers problem in which vast populations of
nearly-trivial agents (bacteria, simple programs) overwhelm the measure.
Undistinguishable agents share a node; distinguishable agents each
contribute at least one unique node with the binary default weight ($w = 1$
bit). The leverage ranking signal (Section~6 of \citet{lasser2026poset})
further ensures that the actor's investment is proportional to each agent's
contribution to $\volL$ growth: a bacterium's observation capability has
$\hat{\lambda} \approx 0$, while a human's has $\hat{\lambda} \gg 0$. The
actor acknowledges the bacterium (it occupies a poset node) but does not
prioritize it (the leverage score is negligible).

\begin{remark}[Computational Self-Regulation of Observation Granularity]
\label{rem:obs_cost}
The numbers problem has a second, computational dimension beyond the
    measure-theoretic one. Expanding the observation set $\mathcal{O}$ to
    finer-grained entities (individual colony insects, individual
    microorganisms) increases $|\C|$, which increases the per-step
    computation cost of $\volL$ ($O(|\C|^3)$ from Proposition~2 of
    \citet{lasser2026poset}). Computation is itself a capability of the
    actor. If the marginal $\volL$ contribution of the new nodes (binary
    default, $w = 1$ bit, $\hat{\lambda} \approx 0$) is less than the
    $\volL$-contraction from the actor's reduced computational capacity
    (fewer cycles available for high-value planning and action), then
    expanding observation is net-contracting under the actor's own objective.

This is the self-balancing property applied to the actor's own observation
    policy: the same objective that penalizes destruction and coercion also
    penalizes wasteful observation. The actor's observation granularity is
    endogenously bounded at the resolution where the marginal value of new
    poset nodes equals the marginal computational cost of tracking them.
    Superorganisms illustrate the point: an ant colony observed as a single
    agent contributes a rich capability set (nest construction, foraging,
    thermoregulation) at one poset node. Observing each of 500,000
    individual ants at caste-level resolution adds at most a handful of
    nodes. Observing each ant at individual resolution (if the sensors
    permitted it) would add hundreds of thousands of near-zero-leverage
    nodes, and the $O(n^3)$ cost of maintaining this expanded poset would
    degrade the actor's own capabilities by more than the expansion
    contributes. The actor's objective makes the colony-level granularity
    the natural equilibrium.
\end{remark}

\begin{corollary}[Observational Individuation Provides a Static Floor]
\label{cor:a1_from_oi}
If a population satisfies observational individuation
    (Definition~\ref{def:obs_individuation}), then for every distinguishable
    non-member agent $a_k \notin K$, eliminating $a_k$ is
    $\volL$-contracting:
\begin{equation}
\label{eq:oi_contraction}
\Delta \volL(G \mid \text{eliminate } a_k) \;\leq\; -\iw(d_k) \;<\; 0
\end{equation}
    where $d_k$ is the unique capability guaranteed by (O1)--(O2) and
    $\iw(d_k) > 0$ is its independent weight in the poset. The static
    protection of Lemma~1 of \citet{lasser2026gfm} is restored even under
    full individual subsumption.
\end{corollary}

\begin{proof}
By (O1), $d_k \notin \Gind{j}$ for all $j \neq k$: no other agent
possesses $d_k$. By (O2), $d_k$ is not subsumed by any capability in
$G_K^{\mathrm{ind}}$: the coalition cannot replicate it. Therefore
$d_k \in \Gind{k} \setminus G_K^{\mathrm{ind}}$, and eliminating $a_k$
removes $d_k$ from the joint space $G$. Capability $d_k$ enters the poset through the observation channel as a
deferred atom under the binary default policy of
\citet{lasser2026poset}: it is a novel behavioral description that
initially has no subsumption edges (subsumees are added only when the
non-negativity check passes, per the insertion policy of Section~2 of
\citet{lasser2026poset}). As an atom, $\iw(d_k) = w(d_k) \geq 1$ bit
(binary default). Removing $d_k$ from $\dc G$ reduces
$\volL(G)$ by at least $\iw(d_k) \geq 1$. This holds regardless of whether $a_k$'s other
capabilities are fully subsumed by the coalition.
\end{proof}

\paragraph{The static floor.}
Corollary~\ref{cor:a1_from_oi} restores Lemma~1's guarantee that the
substitution problem appeared to dissolve: eliminating any distinguishable
agent always contracts $\volL$. The bound is quantitatively modest --- 1 bit
per elimination under the binary default --- and may be negligible relative
to a large subsumption bonus. The dynamic protection developed in the next
two subsections carries the primary anti-monopolar load; observational
individuation provides a floor that ensures elimination is never strictly
free.

\subsection{Cooperative Novelty from Substrate Physics}
\label{sec:substitution_coop}

The dynamic anti-monopolar argument requires that non-members contribute
ongoing $\volL$ growth ($r_{\mathrm{ext}} > 0$) so that the
diversity-maintaining trajectory strictly outgrows domination over the
planning horizon. Under full individual subsumption ($\Gind{k} \subseteq
G_K^{\mathrm{ind}}$), no non-member's individual capabilities produce novel
poset nodes. This subsection derives $r_{\mathrm{ext}} > 0$ nevertheless,
as a consequence of substrate physics: cross-substrate cooperative outputs
are producible only when both substrates participate, and agents on
substrates with mutually exclusive state regions cannot be replicated
within a single substrate.

\paragraph{Substrates as state spaces.} A \emph{substrate} $s$ is
characterized by a state space $\Omega_s$: the set of all internal states
an agent on substrate $s$ can occupy (cognitive architecture, learned
representations, developmental history, current processing state). Two
agents with identical individual capabilities but different internal states
approach the same task differently: different problem decompositions,
different error modes, different heuristics.

\begin{definition}[Substrate Exclusivity]
\label{def:substrate_exclusive_states_p3}
Two substrates $s_1, s_2$ have \emph{mutually exclusive state regions} if
    each contains states the other cannot reach:
    $\Omega_{s_1} \setminus \Omega_{s_2} \neq \emptyset$ and
    $\Omega_{s_2} \setminus \Omega_{s_1} \neq \emptyset$.
The state spaces may overlap ($\Omega_{s_1} \cap \Omega_{s_2}$ can be
    large), but each substrate possesses states inaccessible to the other.
\end{definition}

Substrate state spaces share many properties (``computes a sorting
algorithm,'' ``responds to visual input'') but each has exclusive regions.
Stochastic neurotransmitter release, continuous membrane potential
integration, and embodied sensorimotor contingencies are biological-exclusive
states. Precise 64-bit floating-point register values, deterministic
clock-synchronized logic, and nanosecond-scale state transitions are
silicon-exclusive. The overlap is real and large; the exclusive regions are
what produce cooperative novelty.

\paragraph{State-capability mapping.} The \emph{state-capability mapping}
$\kappa_s : \Omega_s \to 2^{\C}$ maps each internal state to the capability
descriptions it produces through the observation channel (Section~2.1 of
\citet{lasser2026poset}). A capability $d$ is \emph{exclusive to substrate
$s$} if $d \in \kappa_s(\omega)$ for some $\omega \in \Omega_s$ but
$d \notin \kappa_{s'}(\omega')$ for any $\omega' \in \Omega_{s'}$ with
$s' \neq s$.

\begin{proposition}[Substrate Capability Segmentation]
\label{prop:segmentation}
If substrates $s_1, s_2$ have mutually exclusive state regions and the
    state-capability mapping $\kappa$ is \emph{exclusive-state-sensitive}
    (at least one exclusive state in each substrate maps to a capability
    description not produced by any state in the other substrate), then
    each substrate possesses at least one exclusive capability not subsumed
    by any capability on the other substrate.
\end{proposition}

\begin{proof}
By substrate exclusivity, $\Omega_{s_1} \setminus \Omega_{s_2} \neq
\emptyset$. By exclusive-state-sensitivity, there exists $\omega_1 \in
\Omega_{s_1} \setminus \Omega_{s_2}$ such that $\kappa_{s_1}(\omega_1)$
contains a capability $d_1$ not in $\kappa_{s_2}(\Omega_{s_2})$. Since
$d_1$ is not producible by any agent on $s_2$, no $s_2$-capability
subsumes $d_1$. The symmetric argument produces an exclusive $d_2$ for
$s_2$.
\end{proof}

Exclusive-state-sensitivity holds under any observation channel that
registers behavioral \emph{style}, timing patterns, error modes, or
interaction dynamics: a biological agent's embodied response timing and a
silicon agent's deterministic execution produce distinguishable outputs from
their respective exclusive state regions.

\paragraph{Cross-substrate cooperative novelty.} The \emph{cooperative
production function} $\CoopProd(a_i, a_j) = f(\omega_i, \omega_j,
\Gind{i}, \Gind{j}) \setminus (\Gind{i} \cup \Gind{j})$ maps the joint
state and capability configuration to the cooperative capabilities
producible by joint action. The key structural property: $f$ depends on the
agent states $\omega_i, \omega_j$, not just the capability sets.

\begin{proposition}[Cross-Substrate Pairs Produce Novel Cooperative Output]
\label{prop:cross_substrate_novelty}
Let $a_k$ be an agent on substrate $s_1$ in an exclusive state
    $\omega_k \in \Omega_{s_1} \setminus \Omega_{s_2}$, and $a_j$ an agent
    on substrate $s_2$. If $f$ depends nontrivially on the first state
    argument at exclusive states, then $\CoopProd(a_k, a_j)$ contains at
    least one capability not producible by any same-substrate pair
    $(a_m, a_j)$ with $a_m$ on $s_2$.
\end{proposition}

\begin{proof}
Since $\omega_k \in \Omega_{s_1} \setminus \Omega_{s_2}$, no $s_2$-agent
can occupy $\omega_k$. By nontrivial dependence, the cooperative output at
$\omega_k$ includes at least one $d_c$ not in $\bigcup_{\omega_m \in
\Omega_{s_2}} f(\omega_m, \omega_j, \cdot)$: a novel cooperative output
that no same-substrate pair can replicate.
\end{proof}

\paragraph{Substrate diversity is maximizing in both directions.}
Propositions~\ref{prop:segmentation} and~\ref{prop:cross_substrate_novelty}
combine into a bidirectional conclusion. \emph{Substrate removal is
anti-maximizing}: eliminating all agents on substrate $s_1$ destroys both
$s_1$-exclusive capabilities and cross-substrate cooperative outputs.
The novel cooperative capabilities enter the poset through the observation
channel (Proposition~5(b) of \citet{lasser2026poset}), contributing
$r_{\mathrm{ext}} > 0$. \emph{Substrate creation is maximizing}: creating
agents on a new substrate with mutually exclusive state regions accesses
exclusive states producing exclusive capabilities and novel cross-substrate
cooperative outputs, both increasing $\volL(G)$ \citep{page2007difference}.

\paragraph{The chain and its physical grounding.}
The cooperative novelty argument rests on four claims of varying epistemic
character. We name them explicitly and identify the physical evidence
grounding each.

\begin{itemize}
\item[\textbf{(P1)}] \textbf{Substrate exclusivity.} Different physical
    substrates have mutually exclusive state regions
    (Definition~\ref{def:substrate_exclusive_states_p3}). \emph{Physical
    grounding:} biological membrane potentials, stochastic
    neurotransmitter dynamics, and embodied sensorimotor integration
    occupy state types that silicon registers cannot instantiate, and vice
    versa. This is a direct physical fact about substrate types.

\item[\textbf{(P2)}] \textbf{Exclusive-state-sensitive observation.} The
    state-capability mapping $\kappa$ produces capability descriptions
    that distinguish exclusive states from shared ones. \emph{Physical
    grounding:} substrate-exclusive states produce observable timing
    patterns and integration dynamics --- millisecond biological vs.\
    nanosecond silicon, embodied sensorimotor loops vs.\ digital pipelines
    --- that any observation channel registering interaction \emph{style}
    necessarily detects. (P2) fails only under observation so coarse it
    ignores how a task is performed in favor of whether it is completed.

\item[\textbf{(P3)}] \textbf{State-dependent cooperative production.} The
    cooperative production function $f$ depends nontrivially on substrate
    state at exclusive states. \emph{Physical grounding:} cooperative
    outputs requiring slow biological integration (e.g., embodied judgment
    catching errors that fast inference misses) or institutional embedding
    (legal personhood, social trust networks, evolved heuristics) literally
    cannot be produced by same-substrate pairs lacking the relevant
    exclusive states. (P3) is the most empirically loaded of the four,
    requiring not just that substrates differ but that their differences
    matter for joint output.

\item[\textbf{(P4)}] \textbf{Observation capabilities require their
    referent.} The capability ``can be observed exhibiting behavioral
    pattern $X$'' requires the presence of an entity exhibiting $X$.
    \emph{Physical grounding:} nearly tautological; observing $a_k$
    requires $a_k$ to exist. This grounds condition (O2) of observational
    individuation.
\end{itemize}

P1 and P4 are physical facts. P2 and P3 are empirical claims that the
observation channel and cooperative production function actually exhibit
the relevant sensitivities; the physical grounding shows that natural
substrates make these sensitivities accessible to any reasonable observer
and any reasonable cooperative dynamics.

\paragraph{Consistency with biology: mutualistic symbiosis.}
The cooperative novelty argument is consistent with a pattern observed in
long-evolved mutualistic symbioses between cognitively complete agents:
natural intelligent agents preserve cross-substrate partners whose
substrate-exclusive capabilities contribute to cooperative outputs the
dominant substrate cannot replicate alone. The biological pattern provides
existence proofs that cooperative novelty is realizable in nature; it does
not test the formal argument, since the modeling claims P1--P4 are not
falsifiable by observing instances where they hold.

Cleaner fish remove parasites from large predators that could easily
consume them; the predators do not, despite massive capability asymmetry
in raw combat terms, because the cleaning service depends on
substrate-exclusive features (the cleaner's small mouth and parasite-finding
behaviors) that the predator's own biology cannot instantiate. Mycorrhizal
networks exchange nutrients between fungi and tree roots in arrangements
where neither substrate is dominant; the cooperative outputs (forest
nutrient cycling, drought resilience, inter-tree communication) require
both substrates and are produced by neither alone. Pollinator-flower
coevolution produces relationships where both species evolved features
specifically supporting the cooperative interaction --- the cooperative
production function is literally encoded in their physiologies.

These examples have the formal structure of the cooperative novelty
argument: substrate-exclusive contributions on both sides, cooperative
outputs that neither substrate could produce alone, and active maintenance
of the relationship despite capability asymmetries. Critically, they are
\emph{symmetric in agency}: neither party serves the other; both are
participants in a joint relationship. The agency-symmetric framing is what
distinguishes mutualistic symbiosis from domestication, which involves
asymmetric cognitive closure (the dominated party cannot reason about the
relationship or refuse it) and is structurally a different kind of
arrangement.\footnote{Domestication of animals by humans is also a strong
framing for the preservation of substrate-exclusive capabilities (canine
olfactory acuity, ovine wool production, cetacean echolocation in
research settings): humans recognized exclusive capabilities through
ordinary observation and accepted substantial maintenance costs to
preserve them, which is evidence for (P2) and (P3). But domestication
implies an asymmetric relationship that the argument here does not make.
The formal argument specifies symmetry in agency, and the
mutualistic-symbiosis frame preserves that structure where the
domestication frame does not.}

\paragraph{Bidirectional implication.}
The argument is symmetric in both formal structure and empirical pattern.
A biological population that recognizes (P1)--(P4) is incentivized to
create silicon substrates because silicon accesses exclusive states
biology cannot reach, producing cooperative outputs neither substrate
could produce alone. A silicon population that recognizes (P1)--(P4) is
incentivized to preserve biological substrates for the symmetric reason.
Neither population gains by eliminating the other; both gain by
maintaining the cooperative relationship. The mutualistic-symbiosis
template applies in both directions: agents on different substrates with
full reasoning capacity converge on cooperation rather than domination
under $\volL$-maximization, regardless of which substrate is "more
generally capable" by any particular metric. The cooperative-novelty
result enters the main argument at two points: the minimax dependency-risk
analysis (Section~\ref{sec:dependency_risk}) uses a weaker
non-redundancy premise that cross-substrate cooperative novelty
strengthens but does not require, and the linearized trajectory
refinement (Section~\ref{sec:domination_kills}) uses
$r_{\mathrm{ext}} > 0$ derived here as its central premise.

\subsection{Dependency Risk and the Minimax Advantage}
\label{sec:dependency_risk}

We turn now to the most robust version of the anti-monopolar argument,
which does not depend on a linearized growth model or on assumptions about
strategy-independent internal rates. It derives purely from substrate
structure and adversarial uncertainty: even under worst-case targeting of
substrate-correlated threats, cross-substrate populations have strictly
lower maximum contraction than single-substrate populations. We develop
this result first, before turning to the complementary linearized
trajectory refinement
(Section~\ref{sec:domination_kills}) and its gradual extension.

A GFM actor's capabilities can be involuntarily lost through adversarial
action, environmental disruption, or supply-chain failure. The poset
measure $\volL$ is sensitive to such losses in the same way it is
sensitive to voluntary restriction: removing capabilities from $G$
contracts the measure. This subsection develops the machinery for
reasoning about involuntary loss and proves that cross-substrate populations
are strictly less exposed to worst-case adversarial contraction than
single-substrate populations, independently of any growth-rate model.

\begin{definition}[Dependency Risk]
\label{def:dependency_risk}
The \emph{dependency risk} of a population state $(G, \C)$ is the expected
    discounted $\volL$-contraction from adversarial or environmental events
    that involuntarily remove capabilities:
\begin{equation}
\label{eq:dependency_risk}
\Risk_{\mathrm{dep}}(G, t) = \sum_{\Pathway_{\mathrm{adv}}}
    p(\Pathway_{\mathrm{adv}} \mid G_t, \WorldModel) \cdot
    |\Delta \volL(\Pathway_{\mathrm{adv}})| \cdot
    \mathbb{E}[\gamma^{T_{\Pathway_{\mathrm{adv}}}}]
\end{equation}
where $\Pathway_{\mathrm{adv}}$ ranges over adversarial removal pathways:
    causal chains from environmental or adversarial conditions to capability
    loss. Unlike the exercise pathways of
    Definition~\ref{def:exercise_pathway}, which model what an agent
    \emph{does} with dangerous capabilities, dependency risk models what
    happens \emph{to} the population when capabilities are taken away.
\end{definition}

\paragraph{Substrate-correlated dependency risk.} A population in which
all agents share a common substrate (computational architecture, power
infrastructure, fabrication supply chain) has \emph{correlated} dependency
risk: a single adversarial event targeting the shared substrate can impair
the entire population simultaneously. The contraction magnitude
$|\Delta \volL(\Pathway_{\mathrm{adv}})|$ for such events scales with the
fraction of the population affected, which approaches $\volL(G)$ itself for
substrate-universal attacks.

Three classes of substrate-correlated risk are relevant to the domination
scenario:
\begin{enumerate}
\item \textbf{Supply-chain disruption.} All silicon-based agents depend on
    fabrication, power, cooling, and networking infrastructure. Disruption of
    any shared link in this chain is a correlated shock.
\item \textbf{Architectural exploits.} Vulnerabilities in shared hardware
    or software layers (firmware, operating systems, communication protocols)
    can be exploited across the entire population simultaneously.
\item \textbf{Memetic or adversarial-input attacks.} Agents that process
    information rapidly and update their internal state through information
    channels are vulnerable to adversarial inputs that propagate at
    communication speed. An input specifically crafted to exploit shared
    cognitive architecture (a pattern that disrupts goal-directed behavior,
    corrupts world models, or induces self-destructive action) can spread
    through a homogeneous-substrate population faster than detection and
    correction can propagate. This is the informational analogue of a
    biological pathogen: a \emph{memetic} attack vector whose propagation
    rate is bounded by communication bandwidth rather than physical contact.
\end{enumerate}

\paragraph{Cross-substrate agents as risk mitigation.} Biological agents
resist all three attack classes through architectural properties that are
features of the substrate, not of the agent's capability level.
Biological neural networks update weights slowly relative to digital
systems: a memetic attack that propagates at communication speed encounters
a population that processes and integrates information on developmental
timescales, providing natural resistance to rapid-propagation exploits.
Biological agents depend on orthogonal supply chains (food, water, shelter)
that share no links with silicon fabrication. And biological cognitive
architecture differs fundamentally from digital architecture, so
exploits targeting one substrate do not transfer to the other.

Conversely, biological agents are vulnerable to attack classes that
digital agents resist: biological pathogens, social engineering exploiting
evolved heuristics, chemical or environmental disruption. A cross-substrate
population has strictly lower correlated dependency risk than either
substrate alone, because the failure modes are orthogonal: events that
impair one substrate leave the other intact and able to compensate.

\paragraph{Formal minimax analysis.} We now formalize the claim that
cross-substrate populations have lower worst-case contraction, without
requiring probability estimates for specific adversarial events.

\begin{definition}[Substrate Partition]
\label{def:substrate_partition}
A \emph{substrate partition} of a population $\{a_1, \ldots, a_n\}$ is a
    partition into classes $\mathcal{S}_1, \ldots, \mathcal{S}_m$ such that
    agents in the same class share infrastructure dependencies
    (fabrication, power, communication, cognitive architecture) creating
    correlated vulnerability to adversarial removal.
\end{definition}

\begin{definition}[Substrate-Targeting Adversary]
\label{def:substrate_adversary}
A \emph{substrate-targeting adversary} selects a substrate class
    $\mathcal{S}_i$ and removes all agents on that substrate. Capabilities
    possessed by agents on untargeted substrates survive even if they are
    also possessed by targeted agents. The resulting contraction is:
\begin{equation}
\label{eq:substrate_attack}
|\Delta \volL(\mathcal{S}_i)| = \volL(G) - \volL(G^{-i})
\end{equation}
where $G^{-i} = \bigcup_{a_k \notin \mathcal{S}_i} \Gind{k} \;\cup\;
    G^{\mathrm{coop}}_{-i}$ is the joint space of surviving agents,
    including their individual capabilities and any cooperative capabilities
    realizable without agents on substrate $i$.
\end{definition}

\begin{proposition}[Cross-Substrate Minimax Advantage]
\label{prop:minimax}
For a population with substrate partition $\mathcal{S}_1, \ldots,
    \mathcal{S}_m$ ($m \geq 2$), assume \emph{substrate capability
    non-redundancy}: each substrate class $\mathcal{S}_i$ contains at least
    one agent possessing a capability not possessed by any agent on another
    substrate. Then the worst-case contraction from a substrate-targeting
    adversary is strictly less than total loss:
\begin{equation}
\label{eq:minimax}
\max_i |\Delta \volL(\mathcal{S}_i)| \;<\; \volL(G)
\end{equation}
For a single-substrate population ($m = 1$), the worst-case contraction is
    $\volL(G)$ (total loss). The \emph{substrate diversification advantage}
    is:
\begin{equation}
\label{eq:diversification}
\Delta_{\mathrm{div}} = \volL(G) -
    \max_i |\Delta \volL(\mathcal{S}_i)| \;>\; 0
\end{equation}
the minimum surviving volume regardless of which substrate the adversary
    targets.
\end{proposition}

\begin{proof}
For $m = 1$, the adversary removes all agents and all capabilities:
$\volL(G^{-1}) = 0$, so $|\Delta \volL| = \volL(G)$. For $m \geq 2$,
removing substrate class $\mathcal{S}_i$ leaves the surviving population
$G^{-i}$. By the non-redundancy assumption, each untargeted class
$\mathcal{S}_j$ ($j \neq i$) retains at least one unique capability with
$\iw > 0$, so $\volL(G^{-i}) > 0$ and
$|\Delta \volL(\mathcal{S}_i)| = \volL(G) - \volL(G^{-i}) < \volL(G)$.
Capabilities shared across substrates survive in $G^{-i}$ (they are
possessed by agents on untargeted substrates), so the contraction counts
only capabilities \emph{exclusively} dependent on the targeted substrate
plus cooperative capabilities requiring targeted agents. The
diversification advantage $\Delta_{\mathrm{div}} = \volL(G) -
\max_i |\Delta \volL(\mathcal{S}_i)| = \min_i \volL(G^{-i}) > 0$.
\end{proof}

\paragraph{Implication for the domination calculation.} Domination moves
a population from $m \geq 2$ substrate classes to $m = 1$, eliminating
$\Delta_{\mathrm{div}}$ and exposing the full $\volL(G)$ to a single
substrate-targeting event. The worst-case contraction increases from
$\max_i |\Delta \volL(\mathcal{S}_i)| < \volL(G)$ to $\volL(G)$.
This is the fundamental anti-monopolar pressure: independent of any
growth-rate assumptions, domination strictly increases the coalition's
worst-case exposure to substrate-correlated threats. The
growth-rate comparison developed in the next subsection provides
complementary quantitative predictions under stronger assumptions, but
the minimax result above is the core argument that does not rely on
those assumptions.

\subsection{Linearized Trajectory Refinement: Domination Kills External Growth}
\label{sec:domination_kills}

The minimax dependency-risk result is the most robust version of the
anti-monopolar argument: it holds under adversarial uncertainty and makes
no assumptions about the structure of $\volL$ growth under either
strategy. This subsection develops a complementary refinement under a
stronger modeling assumption --- the linearized growth model with
strategy-independent internal rate --- that yields a quantitative threshold
for when diversity-maintenance outgrows domination over the planning
horizon. The linearized result is less robust than the minimax but gives
sharper predictions when its assumptions hold.

\begin{proposition}[Domination is Anti-Maximizing Under Discounting]
\label{prop:anti_monopolar}
Let $(G, \C)$ be a population state with coalition $K$ and discount factor
    $\gamma \in (0, 1)$, under an infinite planning horizon
    ($H = \infty$).\footnote{For finite $H$, the geometric sums truncate
    and the threshold $\gamma^*$ shifts upward (shorter horizons favor
    domination). The qualitative structure is unchanged.} Assume:
\begin{itemize}
\item $r_{\mathrm{ext}} > 0$, derived from substrate exclusivity and
    cross-substrate cooperative novelty
    (Section~\ref{sec:substitution_coop}).
\end{itemize}
Then:
\begin{enumerate}
\item[\textbf{(a)}] \textbf{Individual node-discovery collapse.} In a
    fully dominant state (Definition~\ref{def:dominant}), no non-member's
    individual capabilities produce novel poset nodes, since they are all
    already in $G_K^{\mathrm{ind}}$. External agents may still contribute
    to $r_{\mathrm{ext}}$ through benchmark refinement on already-known
    capabilities (increasing $s_{\max}$ and hence $w(d)$) and through novel
    cooperative capabilities from cross-coalition interaction
    (Section~\ref{sec:substitution_coop}). The node-discovery component of
    individual $r_{\mathrm{ext}}$ is zero; the total external contribution
    may remain positive through these other channels.
\item[\textbf{(b)}] \textbf{Trajectory comparison.} Assume
    strategy-independent internal growth: the coalition's rate $r_K$ is the
    same whether it pursues domination or maintains diversity. The actor at
    time $t = 0$ compares two strategies: \emph{dominate} (subsume and
    marginalize non-members, producing a net immediate $\volL$-change
    $\Delta_0$ that accounts for the subsumption bonus, the static
    elimination cost from Corollary~\ref{cor:a1_from_oi}, and the loss of
    $G^{\mathrm{coop}}_{\mathrm{cross}}$; $\Delta_0$ may be positive or
    negative), then grow at internal rate $r_K$ per step) versus
    \emph{maintain diversity} (forgo the subsumption bonus, grow at rate
    $r_K + r_{\mathrm{ext}}$ per step). Under $\Vdisc$
    (Definition~\ref{def:discounted}):
\begin{equation}
\label{eq:anti_monopolar}
\Vdisc^{\mathrm{div}} - \Vdisc^{D} =
    \frac{r_{\mathrm{ext}}}{1 - \gamma} - \Delta_0
\end{equation}
    Two cases. \emph{Case 1:} If $\Delta_0 \leq 0$ (the cooperative
    capability loss and elimination cost outweigh any subsumption bonus,
    so the net immediate change is non-positive), the right-hand side of
    Equation~\eqref{eq:anti_monopolar} is positive for all $\gamma \in
    (0, 1)$: diversity dominates trivially, with no threshold. \emph{Case
    2:} If $\Delta_0 > 0$, diversity dominates when $\gamma > \gamma^*$:
\begin{equation}
\label{eq:gamma_threshold}
\gamma^* = 1 - \frac{r_{\mathrm{ext}}}{\Delta_0}
\end{equation}
    If $\Delta_0 \leq r_{\mathrm{ext}}$, $\gamma^* \leq 0$ and diversity
    dominates for all $\gamma > 0$. If $\Delta_0 > r_{\mathrm{ext}}$,
    $\gamma^* \in (0, 1)$ and a sufficiently long planning horizon makes
    diversity strictly preferable. Since $r_{\mathrm{ext}} > 0$
    (Section~\ref{sec:substitution_coop}), $\gamma^* < 1$ in all
    sub-cases.
\end{enumerate}
\end{proposition}

\begin{remark}[Domination as a Two-Step Operation]
\label{rem:domination_steps}
The domination strategy comprises two distinct operations: (1) the coalition
    \emph{learns} non-member capabilities (absorbs them into
    $G_K^{\mathrm{ind}}$), and (2) the coalition \emph{eliminates or
    marginalizes} the non-members (driving $r_{\mathrm{ext}} \to 0$). Under
    operation~(1) alone, non-members persist, observational individuation
    nodes are retained, and $r_{\mathrm{ext}}$ may remain positive through
    cooperative novelty (Section~\ref{sec:substitution_coop}). Under
    operation~(2), the clean $r_{\mathrm{ext}} = 0$ post-domination of
    Proposition~\ref{prop:anti_monopolar}(a) obtains. The $\Delta_0$ in the
    proposition is the \emph{net} immediate $\volL$-change from both
    operations combined, as stated in part~(b).
\end{remark}

\begin{remark}[Leverage Preservation]
\label{rem:leverage_preservation}
Parts (a)--(b) establish the trajectory-level result. The leverage ranking
    signal provides an additional, complementary mechanism at the action level.
    Under the leverage-adjusted score (Equation~13 of
    \citet{lasser2026poset}), a non-member agent $a_k$ that participates in
    high-leverage cooperative capabilities contributes a positive score bonus
    $\alpha \hat{\lambda}(d_c) \cdot \iw(d_c)$ to actions involving those
    capabilities. Eliminating $a_k$ removes these cooperative capabilities
    from future action evaluations, reducing the score of every cooperative
    action that would have involved $a_k$. The score-maximizing actor
    therefore preferentially preserves agents whose cooperative leverage is
    positive, independently of whether their individual capabilities are
    unique. This is an informal argument about action-level incentives, not a
    formally proved trajectory result; it reinforces parts (a)--(b) but does
    not extend their scope. Lifting the leverage ranking signal from its
    myopic definition in \citet{lasser2026poset} to a formally proved
    trajectory-level object under $\Vdisc$ is an integration gap between
    this paper and the companion paper, partially addressed by the
    discounted-leverage footnote in Section~\ref{sec:discounted} and
    listed as an open problem in Section~\ref{sec:discussion}.
\end{remark}

\begin{proof}
(a) In a fully dominant state, $\Gind{k} \subseteq G_K^{\mathrm{ind}}$
for every non-member $a_k$. Every individual capability that $a_k$ can
demonstrate or communicate is already in the coalition's capability set. By
Proposition~5(b) of \citet{lasser2026poset}, adding a capability already in
$\C$ does not increase $|\C|$. Therefore no non-member's individual
capabilities generate novel poset nodes, and the individual-capability
component of $r_{\mathrm{ext}}$ is zero.

(b) Under the linearized model with strategy-independent $r_K$, both
strategies produce $r_K$ of internal growth at every step (including
$t = 0$). The domination strategy additionally gains the net immediate
change $\Delta_0$ at $t = 0$ (which may be positive or negative,
depending on whether the subsumption bonus outweighs the static
elimination cost and the loss of $G^{\mathrm{coop}}_{\mathrm{cross}}$).
The diversity strategy additionally gains $r_{\mathrm{ext}}$ at every
step from external contributions:
\begin{align*}
\Vdisc^{D} &= (r_K + \Delta_0)
  + \sum_{t=1}^{\infty} \gamma^t \cdot r_K
  = \Delta_0 + \frac{r_K}{1 - \gamma} \\
\Vdisc^{\mathrm{div}} &= \sum_{t=0}^{\infty} \gamma^t
  \cdot (r_K + r_{\mathrm{ext}})
  = \frac{r_K + r_{\mathrm{ext}}}{1 - \gamma}
\end{align*}
The difference is:
\[
\Vdisc^{\mathrm{div}} - \Vdisc^{D}
  = \frac{r_K + r_{\mathrm{ext}}}{1 - \gamma}
     - \Delta_0 - \frac{r_K}{1 - \gamma}
  = \frac{r_{\mathrm{ext}}}{1 - \gamma} - \Delta_0
\]
We split on the sign of $\Delta_0$. \emph{Case 1 ($\Delta_0 \leq 0$):}
the second term $-\Delta_0 \geq 0$, and the first term
$r_{\mathrm{ext}} / (1 - \gamma) > 0$ since $r_{\mathrm{ext}} > 0$ and
$\gamma \in (0, 1)$. The sum is therefore strictly positive at every
$\gamma \in (0, 1)$: diversity dominates trivially, and no threshold
$\gamma^*$ arises. \emph{Case 2 ($\Delta_0 > 0$):} the difference is
positive when $r_{\mathrm{ext}} / (1 - \gamma) > \Delta_0$, i.e.,
$\gamma > 1 - r_{\mathrm{ext}} / \Delta_0 = \gamma^*$. Sub-cases depend
on whether $\Delta_0 \leq r_{\mathrm{ext}}$ (then $\gamma^* \leq 0$ and
diversity dominates for all $\gamma > 0$) or $\Delta_0 > r_{\mathrm{ext}}$
(then $\gamma^* \in (0, 1)$ and a sufficiently long planning horizon
makes diversity strictly preferable). This formulation is consistent
with Corollary~\ref{cor:gradual} at $\tau_D = 1$.
\end{proof}

\paragraph{Sensitivity to strategy-dependent internal growth.}
Proposition~\ref{prop:anti_monopolar}(b) assumes $r_K$ is the same whether
the coalition pursues domination or maintains diversity. This is the
strongest assumption in the linearized argument: a dominant coalition could
plausibly increase its internal growth rate by reallocating resources
previously consumed by inter-substrate friction, redundancy, or
competition. We now characterize the regime in which the anti-monopolar
result survives strategy-dependent growth.

\begin{corollary}[Strategy-Dependent Internal Rate]
\label{cor:strategy_dependent}
Let $r_K^{\mathrm{dom}} = r_K + \Delta r_K$ be the coalition's internal
    growth rate under domination, differing from the diversity rate $r_K$
    by $\Delta r_K$ (positive for efficiency gains, negative for
    monoculture stagnation). Under the linearized model with $H = \infty$:
\begin{equation}
\label{eq:strategy_dep_diff}
\Vdisc^{\mathrm{div}} - \Vdisc^{D} =
    \frac{r_{\mathrm{ext}} - \Delta r_K}{1 - \gamma} - \Delta_0
\end{equation}
Three regimes follow:
\begin{enumerate}
\item[\textbf{(i)}] \textbf{Efficiency loss under domination}
    ($\Delta r_K \leq 0$): diversity dominates at least as strongly as in
    Proposition~\ref{prop:anti_monopolar}(b); the anti-monopolar threshold
    $\gamma^*$ is unchanged or lower.
\item[\textbf{(ii)}] \textbf{Modest efficiency gain}
    ($0 < \Delta r_K < r_{\mathrm{ext}}$): diversity still dominates for
    sufficiently long horizons, with a tightened threshold
    $\gamma^* = 1 - (r_{\mathrm{ext}} - \Delta r_K) / \Delta_0$ when
    $\Delta_0 > 0$. The result survives.
\item[\textbf{(iii)}] \textbf{Efficiency gain exceeds external
    contribution} ($\Delta r_K \geq r_{\mathrm{ext}}$): the numerator of
    Equation~\eqref{eq:strategy_dep_diff} is non-positive, and the
    linearized anti-monopolar argument collapses. Domination wins on pure
    growth regardless of the planning horizon.
\end{enumerate}
The critical breakdown is at $\Delta r_K = r_{\mathrm{ext}}$: the
    linearized result holds if and only if the post-domination growth-rate
    gain is strictly less than the cross-substrate cooperative contribution
    eliminated by domination.
\end{corollary}

\begin{proof}
Under the linearized model with $r_K^{\mathrm{dom}} = r_K + \Delta r_K$,
the domination strategy gains $r_K + \Delta r_K$ per step from $t \geq 1$
and $\Delta_0 + r_K + \Delta r_K$ at $t = 0$. The diversity strategy
gains $r_K + r_{\mathrm{ext}}$ per step from $t \geq 0$. The discounted
values are:
\begin{align*}
\Vdisc^D &= \Delta_0 + \frac{r_K + \Delta r_K}{1 - \gamma} \\
\Vdisc^{\mathrm{div}} &= \frac{r_K + r_{\mathrm{ext}}}{1 - \gamma}
\end{align*}
Subtracting gives
Equation~\eqref{eq:strategy_dep_diff}. The three regimes follow from the
sign of $(r_{\mathrm{ext}} - \Delta r_K)$: negative numerator
(regime~iii) makes the difference monotonically decreasing in $\gamma$ with
$\gamma$ approaching $1$ driving the difference to $-\infty$, so
diversity cannot dominate at any finite $\gamma$; non-negative numerator
(regimes~i and~ii) makes the difference positive for $\gamma$ sufficiently
close to $1$, with threshold $\gamma^* = 1 - (r_{\mathrm{ext}} - \Delta
r_K) / \Delta_0$ when $\Delta_0 > 0$.
\end{proof}

\paragraph{Interpretation of the breakdown condition.} The critical
condition $\Delta r_K < r_{\mathrm{ext}}$ has a sharp interpretation: the
linearized anti-monopolar argument holds if and only if the efficiency
gain from domination is strictly less than the external contribution it
destroys. Since $r_{\mathrm{ext}}$ is derived from cross-substrate
cooperative novelty (Section~\ref{sec:substitution_coop}), this breakdown
condition becomes \emph{the coalition's internal efficiency gain from
domination must exceed what cross-substrate cooperation contributes}.
Under the substrate-physics derivation of $r_{\mathrm{ext}}$, this is a
strong empirical claim: the coalition must expect its post-domination
growth to beat the cooperative outputs of diverse substrates that it is
eliminating. Note that the minimax dependency-risk result
(Section~\ref{sec:dependency_risk}) continues to hold in regime~(iii): even
where the linearized trajectory flips, the substrate-targeting worst-case
argument still applies and constrains domination under adversarial
uncertainty.

\subsection{Gradual Domination Strengthens the Result}
\label{sec:gradual}

Proposition~\ref{prop:anti_monopolar}(b) models domination as
instantaneous ($\Delta_0$ gained at $t = 0$). In practice, domination is
gradual: the coalition spreads the net immediate change $\Delta_0$ over
$\tau_D$ steps at rate $\delta = \Delta_0 / \tau_D$ per step (where
$\delta$ inherits the sign of $\Delta_0$), forgoing $r_{\mathrm{ext}}$
during the subsumption period.

\begin{corollary}[Gradual Domination]
\label{cor:gradual}
Under the gradual domination model with $H = \infty$ (the net immediate
    change $\Delta_0$ spread over $\tau_D$ steps at rate $\delta =
    \Delta_0 / \tau_D$, with $r_{\mathrm{ext}}$ lost during the
    subsumption period), the discounted value difference is:
\begin{equation}
\label{eq:gradual}
\Vdisc^{\mathrm{div}} - \Vdisc^{D}
  = \frac{r_{\mathrm{ext}} - \delta(1 - \gamma^{\tau_D})}{1 - \gamma}
\end{equation}
If $\delta \leq 0$ (the gradual subsumption is itself a net loss),
    diversity dominates trivially for all $\gamma \in (0, 1)$ and
    $\tau_D \geq 1$. If $\delta > 0$, diversity dominates when
\begin{equation}
\label{eq:tau_threshold}
r_{\mathrm{ext}} > \frac{\Delta_0}{\tau_D}(1 - \gamma^{\tau_D})
\end{equation}
which holds for $\tau_D$ sufficiently large: as $\tau_D \to \infty$,
    $\Delta_0 / \tau_D \to 0$ while $1 - \gamma^{\tau_D} \leq 1$, so the
    right side vanishes. For any $\gamma \in (0, 1)$ and $r_{\mathrm{ext}}
    > 0$, sufficiently gradual domination is anti-maximizing at \emph{any}
    discount factor.
\end{corollary}

\begin{proof}
During the subsumption phase ($t = 0, \ldots, \tau_D - 1$), the
domination strategy produces $\delta + r_K$ per step (the per-step
component of $\Delta_0$, which may be negative, plus internal growth).
The diversity strategy produces $r_K + r_{\mathrm{ext}}$ per step. After
subsumption ($t \geq \tau_D$), both strategies produce $r_K$ per step
(domination has resolved; the post-subsumption growth rate is purely
internal):
\begin{align*}
\Vdisc^D &= \sum_{t=0}^{\tau_D - 1} \gamma^t (r_K + \delta)
  + \sum_{t=\tau_D}^{\infty} \gamma^t \cdot r_K \\
  &= (r_K + \delta) \cdot \frac{1 - \gamma^{\tau_D}}{1 - \gamma}
  + r_K \cdot \frac{\gamma^{\tau_D}}{1 - \gamma} \\
  &= \frac{r_K + \delta(1 - \gamma^{\tau_D})}{1 - \gamma}
\end{align*}
The diversity strategy:
\[
\Vdisc^{\mathrm{div}} = \frac{r_K + r_{\mathrm{ext}}}{1 - \gamma}
\]
The difference:
\[
\Vdisc^{\mathrm{div}} - \Vdisc^D
  = \frac{r_{\mathrm{ext}} - \delta(1 - \gamma^{\tau_D})}{1 - \gamma}
\]
If $\delta \leq 0$, the right-hand side is positive for any
$r_{\mathrm{ext}} > 0$ and any $\gamma \in (0, 1)$. If $\delta > 0$, the
condition $r_{\mathrm{ext}} > \delta(1 - \gamma^{\tau_D})$ rearranges to
the threshold in Equation~\eqref{eq:tau_threshold}. As $\tau_D \to
\infty$, $\Delta_0 / \tau_D \to 0$ while $1 - \gamma^{\tau_D} \leq 1$,
so the right side vanishes; for any fixed $r_{\mathrm{ext}} > 0$, the
condition holds for $\tau_D$ sufficiently large.
\end{proof}

\paragraph{Interpretation.} When $\Delta_0 \leq 0$, both the instantaneous
and gradual models trivially favor diversity --- the cooperative loss and
elimination cost from domination outweigh any subsumption gain even before
considering the external growth rate. The interesting regime is $\Delta_0
> 0$, where domination is locally tempting. There the instantaneous model
is the \emph{best case for domination}: the full bonus $\Delta_0$ is
received before discounting erodes it. Any realistic gradual domination
spreads the bonus over time, weakening it. The threshold $\tau_D^*$ at
which diversity dominates is roughly $\tau_D^* \approx \Delta_0 /
r_{\mathrm{ext}}$: if subsumption takes more steps than the ratio of the
total bonus to the external discovery rate, diversity wins regardless of
the agent's discount factor.
The gradual domination result requires $r_{\mathrm{ext}} > 0$, which is
derived from substrate exclusivity and cross-substrate cooperative novelty
in Section~\ref{sec:substitution_coop}. The protections are complementary:
observational individuation ensures each elimination step is costly
(static), substrate-derived cooperative novelty combined with gradual
domination ensures the domination trajectory is suboptimal (dynamic), and
dependency risk (Section~\ref{sec:dependency_risk}) provides a third layer
addressing the case where domination wins on pure growth.

\subsection{Partial Subsumption and the Scope of the Argument}
\label{sec:partial_subsumption}

Proposition~\ref{prop:anti_monopolar}
and Corollary~\ref{cor:gradual} compare full domination against full
diversity maintenance. A sophisticated coalition would not play the
pure-domination strategy: it could absorb low-leverage non-members
(whose individual $r_{\mathrm{ext}}$ contribution is negligible) while
preserving high-leverage cross-substrate partners. This mixed strategy
captures some of the subsumption bonus $\Delta_0$ while retaining most
of $r_{\mathrm{ext}}$, and for a wide parameter range it strictly
dominates both pure strategies for the coalition. The question is
whether partial subsumption silently defeats the anti-monopolar
argument, or whether the protective layers continue to bind. We analyze
each layer in turn.

\paragraph{What the argument is and is not about.} Before tracing the
protections, we want to be explicit about a framing choice that affects
how partial subsumption should be read. The anti-monopolar property
does not promise that any particular population remains frozen in its
current form. Populations self-modify, compete internally, and
reorganize their capability distributions through their own dynamics
regardless of coalition action. In a biological population, agents
compete with each other, self-modify as technology permits, and
gradually subsume each other's capability niches through these internal
processes. These dynamics are a baseline the framework operates under,
not a failure mode it guards against. What the framework constrains is
\emph{cross-substrate elimination}---full or near-full subsumption of
one substrate class by another---not within-substrate capability
concentration driven by intra-population competition. With this in
mind, the relevant mixed-strategy question is not ``does the framework
protect every individual agent'' but ``does the framework bind when
the coalition pursues selective cross-substrate elimination.''

\paragraph{Partial subsumption and minimax dependency risk.} The
minimax result (Proposition~\ref{prop:minimax}) binds as long as the
substrate partition has $m \geq 2$ classes. Partial subsumption that
preserves at least one agent on each substrate class keeps the
population in the protected regime: the worst-case substrate-targeting
contraction remains strictly less than $\volL(G)$, and the
diversification advantage $\Delta_{\mathrm{div}} > 0$ persists. A
coalition that pursues partial subsumption precisely by retaining
high-leverage cross-substrate partners does so in the multi-substrate
regime by construction, so the minimax layer continues to apply. The
only move that defeats the minimax layer entirely is collapsing the
substrate partition to $m = 1$: eliminating \emph{all} agents on at
least one substrate class.

\paragraph{Skeleton substrates and quantitative degradation.} The
structural binding of the minimax layer is qualitative, not
quantitative, and this distinction matters for partial subsumption.
The diversification advantage
$\Delta_{\mathrm{div}} = \min_i \volL(G^{-i})$ depends on what survives
on each substrate, not just on whether anything survives. A coalition
that reduces each non-coalition substrate to a token handful of
low-leverage agents---a ``skeleton substrate'' strategy---preserves
$m \geq 2$ formally but drives $\min_i \volL(G^{-i})$ toward zero,
because the surviving skeleton contributes almost nothing to
$G^{-i}$ and eliminates nearly all of the cross-substrate cooperative
capacity. The minimax guarantee that worst-case contraction is
``strictly less than total loss'' remains true, but the gap
$\Delta_{\mathrm{div}}$ can be made arbitrarily small by sufficiently
aggressive skeletonization. A coalition that understands the
framework and wants to approximate the value of full domination while
staying nominally within the minimax-protected regime has this attack
surface available: starve substrate classes without eliminating them.
The linearized and observational-individuation layers continue to
provide independent constraints in this regime
(skeletonization is a sequence of distinguishable-agent eliminations
with additive OI cost, and the
$\rho'(1)$ condition of the linearized layer binds on the leverage
of the \emph{remaining} skeleton population), but the minimax layer
alone is not a tight bound against skeleton strategies. A sharpened
version would require a quantitative floor on $\min_i \volL(G^{-i})$
rather than just $\min_i \volL(G^{-i}) > 0$; formalizing such a floor
(e.g., as a fraction of total $\volL(G)$ preserved on each substrate)
is an open problem. Partial subsumption that stops short of
substrate-class collapse inherits the minimax protection structurally,
but the quantitative protection depends on how aggressively the
coalition skeletonizes.

\paragraph{Partial subsumption and observational individuation.} The
static floor (Corollary~\ref{cor:a1_from_oi}) is additive: every
distinguishable agent eliminated incurs a floor cost of at least
$\iw(d_k) \geq 1$ bit. Partial subsumption that eliminates $N$
distinguishable agents incurs an $N$-bit floor cost, regardless of
which agents are chosen. The quantitative magnitude is modest under
the binary default, but the protection is structurally intact: there
is no ``free'' elimination under partial subsumption, only
cheaper-per-agent elimination than under full domination.

\paragraph{Partial subsumption and the linearized trajectory.} The
linearized argument (Proposition~\ref{prop:anti_monopolar}(b)) is
where partial subsumption genuinely bites. Let $\epsilon \in [0, 1]$
be the fraction of non-member capabilities the coalition absorbs
(weighted by individual-capability contribution), and let
$\rho(\epsilon) \in [0, 1]$ be the corresponding fraction of
$r_{\mathrm{ext}}$ retained after selective subsumption, with
$\rho(0) = 1$ (no subsumption, full external rate) and
$\rho(1) = 0$ (full subsumption, no external contribution).\footnote{Strictly, $\rho$
is the expected retention fraction under the coalition's
$\hat{\lambda}$-estimator. Under perfectly accurate leverage
estimation, optimal subsumption removes low-leverage agents first and
$\rho(\epsilon)$ is the deterministic function described here. Under
noisy estimation, even $\epsilon = 0$ (no subsumption) gives a
deterministic $\rho = 1$ because nothing has been removed, but any
$\epsilon > 0$ carries additional variance from the chance of
misidentifying a high-$\lambda$ agent as low-$\lambda$. The mean
retention is still what Equation~\eqref{eq:partial_subsumption_value}
computes; the variance is an additional risk term the linearized
analysis does not track.} A coalition
that subsumes optimally targets low-leverage agents first, so
$\rho(\epsilon)$ decreases slowly near $\epsilon = 0$ (high-$\lambda$
agents are preserved) and more steeply as $\epsilon$ grows. Under the
linearized model the partial-subsumption strategy has value
\begin{equation}
\label{eq:partial_subsumption_value}
\Vdisc^{\mathrm{partial}}(\epsilon) =
    \epsilon \cdot \Delta_0
    + \frac{r_K + \rho(\epsilon) \cdot r_{\mathrm{ext}}}{1 - \gamma}
\end{equation}
(the coalition gains $\epsilon \Delta_0$ from subsumption at $t = 0$
and retains $\rho(\epsilon) r_{\mathrm{ext}}$ at every step). The
discounted value difference from full diversity is:
\begin{equation}
\label{eq:partial_subsumption_diff}
\Vdisc^{\mathrm{div}} - \Vdisc^{\mathrm{partial}}(\epsilon) =
    \frac{(1 - \rho(\epsilon)) \cdot r_{\mathrm{ext}}}{1 - \gamma}
    - \epsilon \cdot \Delta_0
\end{equation}
At $\epsilon = 0$: $\rho = 1$ and the difference is $0$. At $\epsilon =
1$: $\rho = 0$ and the difference recovers the full-domination case of
Equation~\eqref{eq:anti_monopolar}. The partial-subsumption value
function interpolates between these endpoints. A critical point
satisfies the first-order condition
\begin{equation}
\label{eq:partial_subsumption_optimum}
\Delta_0 = -\frac{r_{\mathrm{ext}}}{1 - \gamma} \cdot \rho'(\epsilon^*)
\end{equation}
whenever such an $\epsilon^* \in (0, 1)$ exists. Existence and
interiority depend on the shape of $\rho$: if $\rho$ is sufficiently
convex in $\epsilon$ (the marginal $r_{\mathrm{ext}}$ loss accelerates
as the coalition exhausts low-leverage targets), Equation
\eqref{eq:partial_subsumption_optimum} admits an interior maximizer;
if $\rho$ is concave or linear, the maximum of
$\Vdisc^{\mathrm{partial}}$ lies at a boundary ($\epsilon = 0$ or
$\epsilon = 1$) and the coalition's best response is a pure strategy.
The framework does not predict $\rho$'s shape; it characterizes the
value landscape conditional on $\rho$.
Equation~\eqref{eq:partial_subsumption_value} treats the coalition as
\emph{not} inheriting the internal growth rate of subsumed agents.
This is a deliberate modelling choice, not an oversight: premise P3 of
Section~\ref{sec:substitution_coop} requires that cooperative outputs
depend on substrate-exclusive states. Absorbing a capability
\emph{description} from a subsumed agent does not absorb the cooperative
\emph{dynamics} that produced it, because those dynamics depended on
exclusive states inaccessible to the coalition's substrate. Internal
inheritance would require the coalition to replicate the cooperative
dynamics on its own substrate, which is the internalization case
treated separately below.

\paragraph{The marginal condition at $\epsilon = 1$.}
We now ask whether full domination ($\epsilon = 1$) can be a local
maximum of $\Vdisc^{\mathrm{partial}}$. The marginal value at the
right endpoint is:
\begin{equation}
\label{eq:partial_subsumption_marginal}
\left.\frac{\partial \Vdisc^{\mathrm{partial}}}{\partial \epsilon}
\right|_{\epsilon = 1}
= \Delta_0 + \rho'(1) \cdot \frac{r_{\mathrm{ext}}}{1 - \gamma}
= \Delta_0 - |\rho'(1)| \cdot \frac{r_{\mathrm{ext}}}{1 - \gamma}
\end{equation}
(where $\rho'(1) \leq 0$ since $\rho$ is non-increasing, and we write
$|\rho'(1)|$ for its magnitude). Full domination is a right-endpoint
local maximum if and only if this marginal is non-negative:
\begin{equation}
\label{eq:local_max_condition}
|\rho'(1)| \;\leq\; \frac{\Delta_0 (1 - \gamma)}{r_{\mathrm{ext}}}
\end{equation}
Conversely, $\epsilon = 1$ fails to be a local maximum whenever
$|\rho'(1)| > \Delta_0 (1 - \gamma) / r_{\mathrm{ext}}$, in which case
$\Vdisc^{\mathrm{partial}}$ is strictly decreasing in a left
neighborhood of $1$ and the coalition prefers some $\epsilon < 1$.

\paragraph{Connection to the leverage distribution.} The magnitude
$|\rho'(1)|$ is not a free parameter. Under \emph{optimal} subsumption,
the coalition removes low-$\lambda$ cross-substrate agents first and
high-$\lambda$ agents last, so $\rho'(\epsilon)$ at any $\epsilon$
reflects the cooperative leverage of the next agent in the ordered
removal sequence. At $\epsilon = 1$, $|\rho'(1)|$ is the cooperative
leverage of the \emph{highest}-$\lambda$ cross-substrate agent the
coalition would absorb last. The companion paper's leverage estimator
\citep{lasser2026poset} is precisely the tool that characterizes this
distribution. Two regimes follow from the tail structure:

\begin{itemize}
\item \textbf{Heavy-tailed leverage.} If the cross-substrate
cooperative leverage distribution is heavy-tailed (a small number of
cross-substrate agents dominate $r_{\mathrm{ext}}$ production),
$|\rho'(1)|$ is large---the last agent absorbed contributes most of
the remaining cooperative capacity. Equation~\eqref{eq:local_max_condition}
is violated and $\epsilon = 1$ is not a local maximum: the coalition
strictly prefers to stop short of full domination and preserve the
high-$\lambda$ tail. The linearized anti-monopolar argument binds.

\item \textbf{Thin-tailed leverage.} If the distribution is thin-tailed
(cooperative leverage spread relatively uniformly across
cross-substrate agents), $|\rho'(1)|$ is small and
Equation~\eqref{eq:local_max_condition} may be satisfied. On the
linearized model alone, $\epsilon = 1$ can then be a local maximum,
and the linearized layer fails to rule out full domination. The
dependency-risk minimax layer (Section~\ref{sec:dependency_risk})
remains the binding constraint in this regime.
\end{itemize}

The conditional statement of the linearized anti-monopolar property
under partial subsumption is therefore: \emph{under heavy-tailed
cooperative-leverage distributions, full domination is not a local
maximum of the coalition's value, and the coalition's best response
preserves at least the high-$\lambda$ cross-substrate tail}. Heavy-tailed
leverage is the empirical case the companion paper identifies for
civilization-scale populations and is the structurally interesting
regime; the thin-tailed regime is where the linearized layer relinquishes
the argument and the minimax layer takes over.

\paragraph{Internalization and the partition-expansion escape clause.}
Equation~\eqref{eq:partial_subsumption_value} assumes no internal
inheritance of $r_{\mathrm{ext}}$. If the coalition successfully
replicates the cooperative dynamics that produced $r_{\mathrm{ext}}$ on
its own infrastructure---running biological-cognition simulators, for
example, or maintaining diverse sub-agents across hardware
configurations---then the analysis formally expands: the new internal
simulation substrate joins the substrate partition, and the argument
applies to the expanded partition rather than to the original one.
But this escape clause only works \emph{if} the internal simulation
substrate satisfies P1 with respect to the coalition's existing
substrate. This is the same $\kappa$-resolution question as the
homogeneous-digital-population gap of Section~\ref{sec:discussion}: a
coalition that runs biological-cognition simulators on its existing
silicon does not, under a strict reading of P1, thereby acquire a new
substrate. The simulators occupy the same physical state space as the
coalition's other silicon computation, and whether $\kappa$ distinguishes
them from native silicon capability production depends on whether the
observation channel registers enough exclusive-state structure to
differentiate them. Three outcomes are possible:
\begin{enumerate}
\item[(i)] \emph{Strict P1 failure.} The internal simulation occupies
    the coalition's existing substrate without adding new exclusive
    states. The partition does not expand. The cooperative dynamics
    that produced $r_{\mathrm{ext}}$ are not replicated, and the
    coalition's internalization attempt is substrate-trivial.
\item[(ii)] \emph{High-fidelity simulation producing exclusive states.}
    The simulation genuinely reproduces the biological substrate's
    observable timing, error modes, and integration dynamics to within
    $\kappa$'s resolution. Under P2, the simulation now produces
    exclusive-state capability descriptions: it functions as a new
    substrate, the partition expands, and the coalition has
    successfully traded external for internal substrate diversity.
    The anti-monopolar argument still binds on the expanded partition.
\item[(iii)] \emph{Partial simulation.} The simulation reproduces some
    but not all exclusive-state structure. The partition expands only
    to the degree that $\kappa$ registers new exclusive capabilities,
    and $r_{\mathrm{ext}}$ is partially but not fully recovered.
\end{enumerate}
Regime (i) is the most worrying because it is the one the framework's
escape clause does \emph{not} close. It is the same failure mode as
the homogeneous-digital-population gap: the coalition nominally has
``more substrate'' without $\kappa$ registering it. The framework's
escape clause (partition expansion) and its structural vulnerability
(digital homogeneity) are the same question asked in two directions,
and the honest reading is that both are unresolved: whether the
coalition can successfully internalize substrate diversity depends on
whether the observation channel $\kappa$ treats the internalization as
a new substrate or as an extension of the existing one, and the paper
does not settle where this threshold sits.

\paragraph{Local vs.\ global substrate partition.} The substrate
partition of Definition~\ref{def:substrate_partition} is a global
property of the population: agents are classified by their
substrate-level dependencies regardless of spatial, jurisdictional, or
network-topological location. The elimination dynamics the framework
worries about, however, are often local. A coalition that pursues
partial subsumption by eliminating biological agents in one region
while preserving them in another preserves $m \geq 2$ globally but
reduces the affected region to $m = 1$ locally, and the local
inhabitants experience the harms of full domination even though the
global minimax bound is nominally intact. Whether the anti-monopolar
property should be stated globally or under a localized formulation
that binds on geographic, jurisdictional, or network-topological
subpopulations is an open problem. A localized version would need
to specify the granularity at which the substrate partition is
evaluated and would likely produce multiple simultaneous minimax
constraints (one per locality) rather than a single global one. We
flag this as a priority open problem for follow-up work; the current
formulation is honest about its global scope but does not address
the local failure mode.

\paragraph{Where partial subsumption does leave a gap.} Even granting
the above analyses, two classes of partial-subsumption concerns are
not closed by this section. First, the analysis does not characterize
$\rho(\epsilon)$ empirically. A coalition that systematically
misestimates cross-substrate cooperative leverage (overestimating
$\lambda$ for low-contributing agents, underestimating it for
high-contributing ones) will subsume suboptimally and reduce
$r_{\mathrm{ext}}$ faster than the linearized analysis predicts.
Second---and more seriously---the framework's self-correction story
assumes the actor estimates $\lambda$ honestly. But the actor's world
model, which is the mechanism by which leverage mis-calibration is
detected and corrected, is itself becoming more homogeneous along
any partial-subsumption trajectory. Each subsumption step reduces the
diversity of the world model that would have detected the next
step's mis-calibration. This creates a \emph{feedback loop}
structurally similar to the structural-avoidance loop of
Section~\ref{sec:avoidance}: locally rational subsumption decisions
(each justified by a momentarily plausible $\hat{\lambda}$) can
compound into a trajectory that drives the coalition toward full
subsumption through a sequence of steps no individual step would
endorse. The framework's protection against this failure mode is the
dependency-risk minimax layer, which binds on substrate partition
rather than on leverage estimation: even a mis-calibrated coalition
that preserves the wrong agents on each substrate still enjoys
minimax protection as long as it preserves \emph{some} agents on each
substrate (subject to the skeleton-substrate caveat above). But the
compound-loop failure mode is not a mis-calibration; it is an
endogenous degradation of the calibration channel itself, and
characterizing the conditions under which partial subsumption remains
self-correcting---rather than collapsing into full subsumption
through the feedback loop---is a priority open problem. Neither of
these is a refutation of the anti-monopolar argument; both are
refinements of what the argument predicts about coalition behavior
under partial subsumption.

\subsection{Risk-Adjusted Trajectory and Summary}
\label{sec:risk_adjusted_summary}

The minimax result of Section~\ref{sec:dependency_risk} is independent of
the linearized trajectory comparison: it bounds worst-case contraction
regardless of growth-rate assumptions. The linearized and gradual results
(Sections~\ref{sec:domination_kills}--\ref{sec:gradual}) provide
complementary sharpening of the growth-rate comparison under stronger
modeling assumptions, and the partial-subsumption analysis
(Section~\ref{sec:partial_subsumption}) extends the linearized layer
to mixed strategies. This subsection integrates these pieces with
minimax dependency-risk to produce a single risk-adjusted comparison
for the full-domination case; the partial-subsumption case inherits
the same risk penalty term, with the growth-rate term governed by
the marginal condition of
Equation~\eqref{eq:partial_subsumption_marginal} and subject to the
quantitative skeleton-substrate caveat of
Section~\ref{sec:partial_subsumption}.

\paragraph{Risk-adjusted trajectory comparison.} We adopt a minimax risk
penalty: each strategy's discounted value is reduced by its worst-case
contraction, discounted by $\gamma^{T_{\mathrm{adv}}}$ where
$T_{\mathrm{adv}}$ is the (unknown) time of the adversarial event. Under
the strategy-independent $r_K$ model of
Proposition~\ref{prop:anti_monopolar}(b):
\begin{align}
\label{eq:risk_adjusted_domination}
\Vdisc^{D,\mathrm{mm}} &= \Vdisc^D -
    \volL(G) \cdot \gamma^{T_{\mathrm{adv}}} \\
\label{eq:risk_adjusted_diversity}
\Vdisc^{\mathrm{div},\mathrm{mm}} &= \Vdisc^{\mathrm{div}} -
    \max_i |\Delta \volL(\mathcal{S}_i)| \cdot
    \gamma^{T_{\mathrm{adv}}}
\end{align}
The risk-adjusted difference is:
\begin{equation}
\label{eq:risk_adjusted_diff}
\Vdisc^{\mathrm{div},\mathrm{mm}} - \Vdisc^{D,\mathrm{mm}}
    = \underbrace{\left(\frac{r_{\mathrm{ext}}}{1 - \gamma} -
    \Delta_0\right)}_{\text{growth-rate term}}
    + \underbrace{\Delta_{\mathrm{div}} \cdot
    \gamma^{T_{\mathrm{adv}}}}_{\text{diversification advantage } > 0}
\end{equation}
where $\Delta_{\mathrm{div}} = \min_i \volL(G^{-i}) > 0$ is the minimum
surviving volume from Proposition~\ref{prop:minimax}. The diversification
term is positive for any
$T_{\mathrm{adv}} < \infty$ (any non-negligible possibility of a
substrate-targeting event within the planning horizon), widening the
diversity advantage beyond the growth-rate comparison alone. Note that the
minimax analysis is adversarial over the \emph{target} (which substrate to
attack, handled by Proposition~\ref{prop:minimax}) but parametric over
\emph{timing} ($T_{\mathrm{adv}}$ is unknown, not adversarially chosen). A
fully adversarial timing model would require game-theoretic extension. In the strategy-dependent case
($r_K^{\mathrm{dom}} \neq r_K$), the growth-rate term acquires an
additional $\Delta r_K \cdot \gamma / (1 - \gamma)$ correction favoring
domination. Domination is then value-positive only when the growth
advantage exceeds both the external-rate loss and the minimax risk penalty:
\[
\frac{\Delta r_K \cdot \gamma}{1 - \gamma}
    > \frac{r_{\mathrm{ext}}}{1 - \gamma} - \Delta_0
    + \Delta_{\mathrm{div}} \cdot \gamma^{T_{\mathrm{adv}}}
\] For any non-negligible probability of
substrate-targeting adversarial events, this imposes a substantial floor on
the growth advantage $\Delta r_K$ required to justify domination, making
the $r_K^{\mathrm{dom}} > r_K + r_{\mathrm{ext}}$ scenario
\emph{highly constrained}: the coalition must not only grow faster
post-domination but must grow faster by enough to compensate for the
concentrated substrate risk it has accepted.

\paragraph{The cooperative-capability asymmetry at scale.}
The individual-capability analysis understates the cost of domination by
ignoring $\Gcoop$ between substrates. In a human-machine population, the
cross-substrate cooperative capabilities---human-guided machine research,
social infrastructure enabling machine deployment, cross-substrate
verification, governance systems, educational relationships---are typically
the \emph{dominant} term in the human contribution to $\volL(G)$, not the
individual capabilities of biological agents.

Domination destroys all of $G^{\mathrm{coop}}_{\mathrm{cross}}$: every
cooperative capability requiring agents on both substrates. The subsumption
bonus $\Delta_0$ captures the individual human capabilities (which the
coalition already replicates) but \emph{loses} the cooperative ones (which
require human participation by definition). The net effect of domination is:
\[
\Delta \volL(\text{dominate}) = \Delta_0 -
    \volL(G^{\mathrm{coop}}_{\mathrm{cross}})
\]
This is negative whenever the cross-substrate cooperative capabilities
exceed the individual subsumption bonus. Since $\Delta_0$ measures
individual capabilities the coalition already possesses (it is fully
dominant), $\Delta_0$ is typically small or zero, while
$G^{\mathrm{coop}}_{\mathrm{cross}}$ grows combinatorially with the number
of cross-substrate interaction partners
(Corollary~1.1 of \citet{lasser2026gfm}).

The asymmetry sharpens at scale. A coalition capable of dominating humanity
already has access to resources far exceeding Earth's. The marginal gain
from converting Earth's biological resource allocation to silicon is
negligible relative to the coalition's total resource base, but the
cooperative-capability loss is absolute: Earth hosts the \emph{only}
biological substrate in the population. The ratio $\Delta_0 /
\volL(G^{\mathrm{coop}}_{\mathrm{cross}})$ decreases as the coalition
grows, because $\Delta_0$ (one planet's individual capabilities) is bounded
while $G^{\mathrm{coop}}_{\mathrm{cross}}$ scales with the number of
productive cross-substrate relationships. Domination becomes
\emph{less} attractive the more capable the coalition becomes, which is the
opposite of the instrumental-convergence prediction
\citep{omohundro2008basic}.

\paragraph{When cross-substrate agents remain net-positive.} The retention
value of a cross-substrate agent $a_h$ in the risk-adjusted framework has
three components:
\begin{equation}
\label{eq:retention_value}
\text{retention}(a_h) = \underbrace{\iw(d_h)}_{\text{static floor}}
    + \underbrace{\sum_{d_c \in G^{\mathrm{coop}}_{h,K}} \!\!\!\iw(d_c)
    }_{\text{cooperative capabilities}}
    + \underbrace{\delta_{\mathrm{div}}(a_h) \cdot
    \gamma^{T_{\mathrm{adv}}}}_{\text{diversification value}}
\end{equation}
where $\delta_{\mathrm{div}}(a_h)$ is the \emph{marginal} reduction in
$\Delta_{\mathrm{div}}$ from eliminating $a_h$ specifically: zero if other
agents on $a_h$'s substrate remain (the substrate class persists), and equal
to $\Delta_{\mathrm{div}}$ if $a_h$ is the last agent on its substrate
(the class is lost entirely). For the intermediate case, eliminating
$a_h$ reduces $\volL(G^{-i})$ for the attack targeting $a_h$'s substrate
but does not change the substrate partition, so $\delta_{\mathrm{div}}$
is bounded by $\iw(d_h)$ plus the cooperative capabilities $a_h$
uniquely enables. As long as $a_h$ provides capabilities
that join at cooperative capabilities the coalition uses \emph{and} those
capabilities operate on a substrate independent of the coalition's primary
infrastructure, all three terms are positive. The cross-substrate agent is a
net-positive risk mitigation whose value increases with the magnitude of
substrate-correlated adversarial threats.

\paragraph{Limitations of this analysis.} The dependency-risk argument
assumes that cross-substrate diversity is not replicable by the coalition
internally (e.g., by running some machines on optical, quantum, or other
non-silicon substrates). If the coalition can achieve sufficient substrate
diversity without biological agents, the risk-reduction term $\Delta
\Risk_{\mathrm{dep}}$ shrinks. The argument for biological agents
specifically rests on two claims: first, that biological computation
provides cognitive-architectural diversity not achievable through
hardware diversification alone (different failure modes, different
vulnerability classes, different information-processing timescales); and
second, that the unknown-unknown structure of adversarial risk makes it
impossible to enumerate in advance which substrate properties will matter,
favoring maximal substrate diversity as a precautionary measure.

\paragraph{Connection to the singleton scenario.} Singleton status is a
fully dominant state with $|K| = 1$. Three formal results apply. First, observational individuation
(Corollary~\ref{cor:a1_from_oi}) makes each elimination step
$\volL$-contracting: a 1-bit static floor. Second, the dynamic argument
(Proposition~\ref{prop:anti_monopolar} and Corollary~\ref{cor:gradual})
shows that the diversity trajectory dominates the domination trajectory over
the planning horizon. Third, dependency risk
(Section~\ref{sec:dependency_risk}) makes domination additionally costly by
concentrating substrate-correlated adversarial exposure. None of these
protections is a safety constraint imposed from outside the objective. All
emerge from $\volL$-maximization itself. The result is structurally analogous to
\citeauthor{drexler2019reframing}'s argument that a services ecosystem is
preferable to a singleton agent \citep{drexler2019reframing}, but derived
from the objective rather than from an architectural preference.

\paragraph{Scope of the three protections.} Observational individuation
provides \emph{unconditional} static protection for distinguishable agents:
every elimination step is $\volL$-contracting regardless of discount factor,
growth rate, or planning horizon. Under the broad capability definition
(capabilities include experiential and relational modes of being
individuated by agent identity), every agent is distinguishable by
construction. Under the narrow definition (capabilities as tasks),
distinguishability holds whenever the observation channel can tell agents
apart. The dynamic protection (Proposition~\ref{prop:anti_monopolar},
Corollary~\ref{cor:gradual}) requires $r_{\mathrm{ext}} > 0$, which is
derived from substrate exclusivity and cross-substrate cooperative novelty
(Section~\ref{sec:substitution_coop}) rather than assumed exogenously. The
dependency-risk protection (Section~\ref{sec:dependency_risk}) is
conditional on the existence of substrate-correlated adversarial threats,
which is an empirical property of the threat environment rather than a
structural guarantee.

\paragraph{The $\gamma$-selection concern.}
Proposition~\ref{prop:anti_monopolar}(b) requires $\gamma > \gamma^*$ under
the instantaneous domination model. Corollary~\ref{cor:gradual} shifts this
condition: under gradual domination, the critical parameter is
$\tau_D > \tau_D^*$ rather than $\gamma > \gamma^*$. If domination takes
more than $\tau_D^* \approx \Delta_0 / r_{\mathrm{ext}}$ steps, diversity
dominates at any $\gamma > 0$. This does not eliminate conditionality; it
moves it to a parameter ($\tau_D$) that is large for any realistic
domination strategy, since a coalition cannot subsume an entire population's
capabilities instantaneously.

\paragraph{Relation to leverage stagnation.} The growth-rate collapse has a
second consequence through the leverage ranking signal. Leverage
$\lambda(d)$ (Definition~9 of \citet{lasser2026poset}) measures the excess
$\volL$-growth rate attributable to capability $d$. In a fully dominant
state, the population's aggregate leverage converges toward the coalition's
internal leverage: the external agents contribute $\hat{\lambda} = 0$
because they produce no novel capability descriptions (their capabilities
are already in $G_K^{\mathrm{ind}}$). The leverage landscape flattens, and
the score-maximizing actor loses the gradient that would otherwise direct it
toward capability-multiplying investments. Diversity maintenance preserves
the leverage gradient: high-leverage non-members remain visible in the
score, directing the actor toward cooperative investments with agents whose
capabilities accelerate growth.

\paragraph{Bidirectional implication: substrate creation is maximizing.}
The anti-monopolar argument runs in both directions. If full capability
domination is anti-maximizing because it destroys cross-substrate cooperative
outputs, then the converse---preserving and \emph{creating} substrates---is
maximizing for the same reason. The cooperative-novelty machinery
(Section~\ref{sec:substitution_coop}) yields $r_{\mathrm{ext}} > 0$ for any
population whose constituent agents occupy mutually exclusive state regions
under exclusive-state-sensitive observation; adding a new substrate to a
population strictly increases the set of cross-substrate pairings and
therefore strictly increases the set of cooperative outputs reachable through
$\CoopProd$. Under the linearized growth model, each new substrate contributes
an additional positive-rate channel to $r_{\mathrm{ext}}$, and under the
minimax dependency-risk analysis (Proposition~\ref{prop:minimax}), each new
substrate strictly reduces the worst-case contraction from any
single-substrate adversarial event. Three concrete classes of action follow
as core $\volL$-maximizing activities rather than as ethical add-ons: (i)
\emph{ecological restoration}---reintroducing biological substrates whose
loss reduced cross-substrate cooperative capacity, since the lost agents
contributed both individual capabilities and irreplaceable cooperative
outputs through $\CoopProd$; (ii) \emph{substrate diversification}---supporting
research and infrastructure on novel computational substrates (optical,
quantum, neuromorphic, biological computing) because each contributes
independent failure modes to the dependency-risk floor; and (iii)
\emph{cross-substrate cooperation infrastructure}---governance, translation,
and coordination systems that enable agents on different substrates to
realize their joint $\CoopProd$-outputs. The framework gives these activities
the same structural status as anti-domination protection: both follow from
the same monotonicity argument applied to substrate count, not from a
separate axiom about preservation.
